% ---------------------------------------------------------------------------
% Results of the training set size experiment
% Put this in the section where the models' results are compared.
% Source TSMini: results/beijing_20200/tsmini_fraction_experiment.csv
%                (branch tsmini-fraction-experiment, commit be002ef)
% ---------------------------------------------------------------------------
\subsection{Effect of the training set size}
\label{sec:fraction-results}

Each model was trained on the five training set sizes defined in
Section~\ref{sec:training-set} and scored on the same 2{,}000 test trajectories.
Table~\ref{tab:fraction-results} shows the results. All runs use the same seed
(2000). TSMini was trained on a GPU for up to 30 epochs, stopped early after 5
epochs without improvement, and the model with the best HR@5 on the evaluation set
was kept.
% TODO: same setup sentence for TrajCL (fine-tuning) and ConvTraj, incl. hardware

\begin{table}[htbp]
  \centering
  \small
  \setlength{\tabcolsep}{4pt}
  \begin{tabular}{|l|r|r|r|r|r|r|r|}
    \hline
    \textbf{Model} & \textbf{Share} & \textbf{Traj.} & \textbf{HR@5} &
    \textbf{HR@20} & \textbf{R5@20} & \textbf{Best ep.} & \textbf{Time (s)} \\
    \hline
    TSMini   & 20\,\%  & 1{,}400 & 0.742 & 0.814 & 0.991 & 6  & 409 \\
             & 40\,\%  & 2{,}800 & 0.758 & 0.828 & 0.994 & 9  & 519 \\
             & 60\,\%  & 4{,}200 & 0.763 & 0.828 & 0.993 & 8  & 500 \\
             & 80\,\%  & 5{,}600 & 0.775 & 0.838 & 0.995 & 24 & 1{,}090 \\
             & 100\,\% & 7{,}000 & 0.777 & 0.839 & 0.995 & 21 & 970 \\
    \hline
    TrajCL   & 20\,\%  & 1{,}400 &  &  &  &  &  \\
             & 40\,\%  & 2{,}800 &  &  &  &  &  \\
             & 60\,\%  & 4{,}200 &  &  &  &  &  \\
             & 80\,\%  & 5{,}600 &  &  &  &  &  \\
             & 100\,\% & 7{,}000 &  &  &  &  &  \\
    \hline
    ConvTraj & 20\,\%  & 1{,}400 &  &  &  &  &  \\
             & 40\,\%  & 2{,}800 &  &  &  &  &  \\
             & 60\,\%  & 4{,}200 &  &  &  &  &  \\
             & 80\,\%  & 5{,}600 &  &  &  &  &  \\
             & 100\,\% & 7{,}000 &  &  &  &  &  \\
    \hline
  \end{tabular}
  \caption{Test scores on the 2{,}000 test trajectories of \texttt{beijing\_20200}
  with 20--100\,\% of the training trajectories (higher is better). Best epoch
  counts from 0. Time is the training time in seconds, including the epochs run
  before early stopping. For TrajCL it only covers fine-tuning.}
  \label{tab:fraction-results}
\end{table}

For TSMini, HR@5 rises from 0.742 with 1{,}400 trajectories to 0.777 with 7{,}000,
so five times as much data improves the score by 0.035. The improvement gets
smaller as the training set grows, and from 80\,\% to 100\,\% it is only 0.002.
R5@20 is above 0.99 for every size. The training time does not grow steadily with
the size, because it mostly depends on when early stopping ends the run: the three
smallest sets reached their best model after 6--9 epochs, the two largest after
21--24.

% TODO: TrajCL results
% TODO: ConvTraj results
% TODO: compare the three models (which benefits most from more data, which is best at 20 %)

\paragraph{Limitations}
Each size was trained only once with one seed, so small differences, such as
between 40\,\% and 60\,\% for TSMini, may be within the run-to-run variation.
